\section{Correction of the cone models}
\label{cc:section}

Throughout this section,
\[
 S=C_0((0,1])\otimes P,\qquad D=Q_\infty.
\]
We first correct the multiplicative defects of the map in
\Cref{tc:rough-model}.  The correction takes place on a hereditary
support containing the defects and a small cone representation.
Its ideal estimates are stronger than equivariance for the projected
labels.  This strength is what permits the cone homotopy to be used
with the finer labels of $\Prim(S)$.

\subsection{The error support and its ideal estimates}

\begin{lemma}[A hereditary support for the defects]\label{cc:defect-support}
Let $h:S\to D$ be \cpc and suppose that, for every open $U\subset X$,
\begin{equation}\label{cc:mixed-defects}
 h(sz)-h(s)h(z),\quad h(zs)-h(z)h(s)\in J(U)
 \qquad(s\in S(U),\ z\in S).
\end{equation}
Let $E=C^*(h(S))$ and let $I\lhd E$ be the ideal generated by
$h(ab)-h(a)h(b)$, for $a,b\in S$.  Then
\begin{equation}\label{cc:annihilation}
 h(S(U))I+Ih(S(U))\subset J(U).
\end{equation}
There is a locally small cone homomorphism $\theta:S\to D$ such that,
after placing $h$ and $\theta$ in separate matrix corners, a
$\sigma$-unital hereditary algebra $H\subset M_2(D)$ has the following
properties.  The first-corner map $h$ idealizes $H$, its multiplicative
defects belong to $H$, and $\theta(S)\subset H$.  Moreover, for every
$x\in S_+$,
\begin{equation}\label{cc:strong-support}
 h(S(U_x))H+Hh(S(U_x))\subset H_x,
 \qquad H_x:=\Ideal_H(\theta(x)).
\end{equation}
Here $\theta(S(U))\subset J(U)$, with the matrix conventions of
\Cref{st:folding}.
\end{lemma}

\begin{proof}
Fix $U$ and pass to $D/J(U)$, writing bars for the quotient maps.
For $s\in S(U)$ and $a,b\in S$, the mixed multiplication identities give
\[
 \bar h(s)\bigl(\bar h(ab)-\bar h(a)\bar h(b)\bigr)
 =\bar h(sab)-\bar h(sa)\bar h(b)=0.
\]
The analogous right product is zero.  More generally a word in
$\bar h(S)$ between $\bar h(s)$ and a defect can be absorbed into
the localized input one factor at a time.  The resulting input is
still in $S(U)$.  Polynomial approximation in $E$ therefore proves
\eqref{cc:annihilation} for the whole generated ideal $I$.

The algebra $E$, and hence $I$, is separable.  Apply
\Cref{sm:coefficient-domination} to the bounded bilinear map
$(s,b)\mapsto h(s)b$, with $b\in I$.  Include its adjoints in the
same countable family.  This gives a locally small cone homomorphism
$\theta$ satisfying, simultaneously for every $x\in S_+$,
\begin{equation}\label{cc:ambient-domination}
 h(S(U_x))I+Ih(S(U_x))
 \subset\Ideal_{M_2(D)}(\theta(x)).
\end{equation}
The first-corner and second-corner placements are understood in this
formula; the ideal on the right is the ambient matrix ideal.

Choose strictly positive elements $e\in I_+$ and
$k\in\theta(S)_+$, omitting a zero summand when necessary, and set
\begin{equation}\label{cc:H-definition}
 H=\Her_{M_2(D)}(e\oplus k).
\end{equation}
It is $\sigma$-unital, since $e\oplus k$ is strictly positive for $H$.
No separability assertion about $H$ is needed.  Since $I\lhd E$,
$h(a)I+Ih(a)\subset I$ for $a\in S$.  Functional-calculus
approximate units of $I$ show that the first-corner map idealizes
the first support.  They also show that it preserves all the
off-diagonal corners in \eqref{cc:H-definition}: on the dense
linear span of products
\[
 f(e\oplus k)\,z\,g(e\oplus k),\qquad z\in M_2(D),
\]
left multiplication by $h(a)$ replaces the first support factor by
an element of $I$, and right multiplication has the corresponding
property.  Approximation by the support approximate unit proves
idealization of $H$.  The defects themselves lie in $I\subset H$.

For $a\in S(U_x)$, the same calculation, now using
\eqref{cc:ambient-domination}, puts $h(a)H$ and $Hh(a)$ in
$H\cap\Ideal_{M_2(D)}(\theta(x))$.  The hereditary ideal
correspondence gives
\[
 H\cap\Ideal_{M_2(D)}(\theta(x))
 =\Ideal_H(\theta(x)).
\]
This proves \eqref{cc:strong-support}; equivalently, one may apply
\Cref{un:hereditary-support}.  All matrix placements preserve the
local smallness of $\theta$.
\end{proof}

We use the right Hilbert $H$-module convention.  Write
$\mathcal L_H(\mathcal E)$ and $\mathcal K_H(\mathcal E)$ for
adjointable and compact operators on a Hilbert $H$-module $\mathcal E$.
For an ideal $L\lhd H$, let $\mathcal E L$ denote the closed span
of $\{\xi b:\xi\in\mathcal E, b\in L\}$.

\begin{lemma}[A dilation compatible with fine ideal labels]\label{cc:scaled-dilation}
Let $d:S\to\mathcal L_H(H)=M(H)$ be the multiplier map induced by
$h$ in \Cref{cc:defect-support}.  There is a countably generated
Hilbert $H$-module $\mathcal E'$ and a representation
\begin{equation}\label{cc:dilation}
 \Pi:S\longrightarrow\mathcal L_H(H\oplus\mathcal E'),
 \qquad
 \Pi(a)=
 \begin{pmatrix}
 d(a)&C(a)\\ B(a)&d'(a)
 \end{pmatrix},
\end{equation}
whose off-diagonal entries are compact.  The maps $d,d'$ are c.p.c.
The module may be stabilized, retaining the first copy of $H$, so
that $\mathcal E'$ is the standard countable module.

Give $S$ its tight $\Prim(S)$-labels and give $H$ the order-preserving
assignment
\begin{equation}\label{cc:fine-labels}
 H(O)=\Ideal_H(\theta(S(O)))
 \qquad(O\subset\Prim(S)\text{ open}).
\end{equation}
Let $\alpha_t$ be cone scaling,
\[
 (\alpha_t a)(r)=a(tr),\qquad 0\le t\le1,
\]
where $a(0)=0$.  Every $\Pi_t:=\Pi\alpha_t$ is weakly
$\Prim(S)$-nuclear for \eqref{cc:fine-labels}.  The same is true
of its diagonal compressions $d_t,d'_t$.
\end{lemma}

\begin{proof}
Use the Hilbert-module Stinespring construction in
\cite[proof of Theorem~3, pp.~140--141]{Kasparov}, applied to the
unital extension $d^+$ of $d$ on the forced unitization of $S$.
The source is separable and $H$ has a countable approximate unit,
as required there.
The resulting representation acts on the completion of the
algebraic Stinespring module.  Its adjointable isometry
$V:H\to\mathcal E$ is given by
$V\xi=1\otimes\xi$ and
$V^*(b\otimes\eta)=d^+(b)\eta$, so
$V^*\Pi(a)V=d(a)$.  The module is countably generated:
use a countable dense subset of the unitized source and a countable
approximate unit of $H$.  Decompose
$\mathcal E=VH\oplus\mathcal E'$ and identify $VH$ with $H$.

For $B(a)=(1-VV^*)\Pi(a)V$,
\begin{equation}\label{cc:compact-offdiagonal}
 B(a)^*B(a)=d(a^*a)-d(a^*)d(a)\in H
 =\mathcal K_H(H).
\end{equation}
An adjointable operator whose square modulus is compact is compact,
by functional calculus in the compact-operator ideal.  Thus $B(a)$
is compact, and so is $C(a)=B(a^*)^*$.  Kasparov stabilization
\cite[Theorem~2]{Kasparov} of $\mathcal E'$ after adding a zero
representation identifies its
enlargement with the standard module.  This leaves $VH$ fixed and
preserves compactness and the indicated diagonal compressions.

We check the ideal assertion before using the homotopy.  For
$a\in S(U_x)$, coefficients on the dense Stinespring generators
have the form
\begin{equation}\label{cc:Stinespring-coefficients}
 \langle b\otimes\xi,\Pi(a)(c\otimes\eta)\rangle
 =\xi^*d(b^*ac)\eta\in H_x,
\end{equation}
where $b,c$ belong to the forced unitization.  The containment uses
$b^*ac\in S(U_x)$ and \eqref{cc:strong-support}.  It extends to
all vectors.  Passing to the quotient module by $\mathcal E H_x$
shows that $\Pi(a)\mathcal E\subset\mathcal E H_x$; the same
holds for the adjoint.  Stabilization and module unitaries preserve
these ideal submodules.

Now take $a\in S(O)$ and put $x=|a|$.  Scaling does not enlarge its
\emph{projected} support:
\begin{equation}\label{cc:projected-scaling}
 \alpha_t(a)\in S(U_x)\quad(0\le t\le1).
\end{equation}
Indeed a point quotient in the $P$ coordinate which kills $a$
also kills every rescaling of $a$.  Since $x\in S(O)$,
$H_x\subset H(O)$, and \eqref{cc:Stinespring-coefficients}
applied to $\alpha_t(a)$ proves weak equivariance of $\Pi_t$.
It proves the same assertion for the compressions.  Source ideals
and quotients are nuclear, so these coefficient conditions also
give weak nuclearity, as in \Cref{un:absorption}.

This argument does not assert that $\alpha_t$ preserves $S(O)$.
The stronger estimate for the entire projected ideal
$S(U_x)$ is the reason that the scaled representation has the
required fine ideal bounds.
\end{proof}

For later use, if $U\subset X$ and $O_U=(0,1]\times U$, then
\begin{equation}\label{cc:local-compact-ideal}
 H(O_U)\subset H\cap M_2(J(U)).
\end{equation}
This follows from the ideal property of $J(U)$ and the local
smallness of $\theta$.  If a compact-valued map is the difference
of two weakly equivariant maps, its value on $S(O_U)$ belongs to
the compact operators over $\mathcal E H(O_U)$.  To see this,
pass to the quotient module by $\mathcal E H(O_U)$: both maps
vanish there on this input, and the kernel of the quotient map on
compact operators is precisely
$\mathcal K_H(\mathcal E H(O_U))$.  Thus every such compact
correction is locally in $J(U)$.

\subsection{A finite correction on the original module}

The dilation now respects every fine ideal label, including along
cone scaling.  We will compare nearby times in a finite subdivision of
that homotopy.  Absorption makes a finite chain of dilations available as
a compact perturbation of a repeat of $\theta$.  Shifting the first
summands along this chain pairs each complementary summand with a nearby
compression; the unpaired summand lies near the zero endpoint.  The
resulting map is close to a representation.

We then truncate only the compact perturbation.  This produces a finite
matrix correction with the original $h$ still in its first corner, so
its square moments and lower ideal support remain available.

\begin{lemma}[A finite correction retaining the original corner]\label{cc:finite-correction}
Under the hypotheses of \Cref{cc:defect-support}, let
$\mathcal F$ be a finite subset of the unit ball of $S$ and let
$\eta>0$.  There are an integer $N$, a bounded linear
$*$-preserving map $L:S\to M_{N+1}(M_2(D))$, and a bounded linear
map $K_N:S\to M_{N+1}(H)$ such that
\begin{equation}\label{cc:finite-form}
 L(a)=h(a)\oplus\theta(a)^{\oplus N}+K_N(a),
\end{equation}
and, for $a,b\in\mathcal F$,
\begin{equation}\label{cc:finite-estimates}
 \|L(a)\|\le\|a\|+\eta,
 \qquad
 \|L(ab)-L(a)L(b)\|\le3\eta+2\eta^2.
\end{equation}
For every open $U\subset X$,
\begin{equation}\label{cc:finite-locality}
 K_N(S(U))\subset M_{N+1}(H\cap M_2(J(U))).
\end{equation}
The first summand in \eqref{cc:finite-form} is the original $h$,
with only zero corner padding.
\end{lemma}

\begin{proof}
\smallskip\noindent\emph{Absorb a finite chain of dilations.} 
Include $\mathcal F^*$ and all products from $\mathcal F$ in the
finite list on which estimates will be imposed.  By norm continuity
of $t\mapsto\alpha_t(a)$, choose
\[
 0=t_0<t_1<\cdots<t_k=1
\]
so that adjacent values differ in norm by less than $\eta$ on
this enlarged list.  Write $d_i=d_{t_i}$, $d'_i=d'_{t_i}$ and
$\Pi_i=\Pi_{t_i}$, using the same first module $H$ and the same
complementary module in all these dilations.

For the absorption application put
$\widehat\theta(a)=\theta(a)\otimes e_{11}\in H\otimes\K$,
with target labels $H(O)\otimes\K$.
If $O_a$ is the fine support of $a\in S_+$, then
$S(O_a)=\Ideal_S(a)$ and
\[
 \Ideal_H(\theta(S(O_a)))=\Ideal_H(\theta(a)).
\]
Thus $\widehat\theta$ is full for these fine labels, independently
of any coarse fullness assertion.  Apply \Cref{un:absorption}
to a literal stable infinite repeat $\Theta$ of
$\widehat\theta$ and the finite sum of the $\Pi_i$.
On the standard Hilbert $H$-module, this repeat can be reindexed as
\begin{equation}\label{cc:stable-repeat}
 \Theta(a)=\operatorname{diag}
       (\theta(a),0,\theta(a),0,\ldots).
\end{equation}
Indeed each repeated copy of $\widehat\theta$ acts by $\theta$
on one $H$ slot and by zero on the other slots of its standard
module.  Both the active slots and the remaining zero slots are
countably infinite, giving \eqref{cc:stable-repeat} by a fixed
permutation.  We retain every zero slot throughout absorption and
finite compression.  The bare diagonal repeat with all these zero
slots removed need not absorb a zero representation.
The hypotheses hold: the tightly labeled source is separable,
exact and lower semicontinuous; $H\otimes\K$ is $\sigma$-unital
and stable; and \Cref{cc:scaled-dilation} supplies weak nuclearity
of the representations being absorbed.  Absorption identifies
the repeat module with a module carrying
\[
 \Theta\oplus\Pi_1\oplus\cdots\oplus\Pi_k
\]
so that its difference from the original repeat is compact on
every input.  More precisely, absorption in
\cite[Definition~5.4]{Gabe} gives a repeat-side module unitary $U$
with
\[
 U^*(\Theta\oplus\textstyle\bigoplus_i\Pi_i)(a)U-\Theta(a)
 \in\mathcal K_H(H^{\oplus\infty})\qquad(a\in S);
\]
the compact membership holds for every input for each chosen
unitary, in addition to the finite-test norm approximation.
Only this compact-membership conclusion from absorption is needed
below.  Its finite-test approximation error does not enter our bound:
after the unitary identification we keep the absorbed representation
itself.  All these identifications act on the repeat side; the initial
module carrying $d=d_k$ is fixed.

\smallskip\noindent\emph{Shift the first summands along the chain.} 
Denote the initial copy of
$H$ by $H_0$, the first copy in $\Pi_i$ by $H_i$, and its
complement by $\mathcal E'_i$.  Delete the compact off-diagonals
of the absorbed $\Pi_i$.  Apart from the unused repeat, the
diagonal summands are
\begin{equation}\label{cc:before-regrouping}
 d_k\big|_{H_0}\ \oplus
 \bigoplus_{i=1}^k
 \left(d_i\big|_{H_i}\oplus d'_i\big|_{\mathcal E'_i}\right).
\end{equation}
Pair $H_0$ with $\mathcal E'_k$, and, for $i<k$, pair
$H_{i+1}$ with $\mathcal E'_i$.  The unused first module is
$H_1$.  Restore the compact off-diagonals of $\Pi_i$ in the pair
whose second component is $\mathcal E'_i$.  In this ordered
decomposition the resulting map is
\begin{equation}\label{cc:regrouped-map}
 T'=\Theta\oplus d_1\oplus
 \bigoplus_{i=1}^{k-1}
 \left[\Pi_i+
 \begin{pmatrix}d_{i+1}-d_i&0\\0&0\end{pmatrix}\right]
 \oplus\Pi_k.
\end{equation}
It is a compact perturbation of the map in
\eqref{cc:before-regrouping}, including its repeat: the diagonal
entries were only permuted, and the restored entries are compact.
The representation
\begin{equation}\label{cc:comparison-representation}
 R'=\Theta\oplus0\oplus
 \bigoplus_{i=1}^{k-1}\Pi_i\oplus\Pi_k
\end{equation}
is within $\eta$ of $T'$ on the enlarged finite list.  Indeed
$\|d_1(a)\|\le\|\alpha_{t_1}(a)\|<\eta$ there, since
$\alpha_0=0$, and every other difference in
\eqref{cc:regrouped-map} has norm less than $\eta$.
No compactness of $d_{i+1}(a)-d_i(a)$ has been used.

Pull both maps back through the block permutation and the
absorption identification to the original module
\[
 \mathcal L=H\oplus H^{\oplus\infty}.
\]
Writing $T_0=d\oplus\Theta$, we obtain
\begin{equation}\label{cc:T-plus-K}
 T=T_0+K,\qquad
 K(a)\in\mathcal K_H(\mathcal L),\qquad
 \|T(a)-R(a)\|<\eta
\end{equation}
on the enlarged list, where $R$ is a representation.
The maps in this construction are bounded linear and
$*$-preserving.  They are weakly equivariant: the block
compressions and off-diagonal entries have the ideal bounds of
\Cref{cc:scaled-dilation}, and any Hilbert-module unitary preserves
the submodule associated to an ideal.  Thus
\eqref{cc:local-compact-ideal} applies to $K$.
For $a,b\in\mathcal F$, comparison with the contractive
representation $R$ gives
\begin{equation}\label{cc:T-estimates}
 \|T(a)\|\le\|a\|+\eta,
 \qquad
 \|T(ab)-T(a)T(b)\|\le3\eta+\eta^2.
\end{equation}

\smallskip\noindent\emph{Compress the compact correction.} 
Let $E_N$ be the projection onto the initial copy of $H$ and the
first $N$ slots in the \emph{original} diagonal repeat
\eqref{cc:stable-repeat}, including both active and zero slots.  It
commutes with $T_0(a)$ for all $a$.  The projections $E_N$ tend
strictly to the identity on the compact operators.  Choose $N$
so that both tails of every $K(a)$ on the enlarged list have
norm less than $\eta$.  Then
\begin{align}
 &E_NT(ab)E_N-E_NT(a)E_NT(b)E_N\notag\\
 &\quad=E_N\bigl(T(ab)-T(a)T(b)\bigr)E_N
       +E_NK(a)(1-E_N)K(b)E_N.
 \label{cc:compression-defect}
\end{align}
The second term has norm at most $\eta^2$.  Compression preserves
the norm estimate in \eqref{cc:T-estimates}.

The entries of $E_NK(a)E_N$ lie in $H$, so the compressed map
has an actual ambient realization with diagonal baseline $h$
followed by the retained $\theta$ and zero slots.  To obtain
exactly \eqref{cc:finite-form}, let $Z_N(a)$ be the finite diagonal
matrix which places $\theta(a)$ in each retained zero repeat slot
and zero in every other slot, including the initial one.  Define
\begin{equation}\label{cc:zero-slot-rebasing}
 K_N(a)=E_NK(a)E_N-Z_N(a).
\end{equation}
Replacing each zero baseline slot by $\theta(a)$ and subtracting
it in \eqref{cc:zero-slot-rebasing} leaves the compressed map $L$
unchanged.  The map $Z_N$ is bounded linear and $*$-preserving,
has entries in $H$, and is locally in $J(U)$ because $\theta$ is
locally small.  Thus $K_N$ has the required compact and local
properties.  This rebasing is performed only after finite
compression; the tail estimate used the original $T_0$ with
all its zero slots.

\smallskip\noindent\emph{Return to the ambient sequence algebra.} 
To transfer the multiplier estimates, let $p_H$ be the support
projection of $H$ in $M_2(D)^{**}$.  Idealization implies that
$h(a)$ commutes with $p_H$.  The correction entries, and the
values of $\theta$, are supported on $p_H$.  On its matrix
amplification the ambient map is exactly the compressed
multiplier map just estimated.  On the orthogonal complement
all corrections vanish; the remaining first-corner map $h$
is contractive and multiplicative there, since its defects lie
in $I\subset H$.  Taking the maximum of these two reducing
parts proves \eqref{cc:finite-estimates} in the ambient norm.
Finally \eqref{cc:finite-locality} follows from the local
compact ideal statement for $K$.
\end{proof}

\subsection{Passage to a corrected sequence homomorphism}

\begin{proposition}[Cone correction]\label{cc:correction}
Suppose $h:S\to D$ is \cpc and satisfies
\eqref{cc:mixed-defects}.  After trace-preserving coordinate
corner padding, there is a homomorphism $\lambda:S\to D$ such
that
\begin{equation}\label{cc:corrected-difference}
 (\lambda-h^{\mathrm{pad}})(S(U))\subset J(U)
 \qquad(U\subset X\text{ open}).
\end{equation}
The padding can be chosen to fix the constant copy of $Q$ in
the sequence quotient.  Consequently it preserves the constant
ideal comparisons and the regularized trace formulas prescribed
for $h$.
\end{proposition}

\begin{proof}
If the defect ideal is zero, $h$ is already a homomorphism.
Otherwise use the preceding construction.  Nuclearity and separability
of $S$ provide a \cpc coordinate lift of $h$ by
\cite[Theorem~3.10]{ChoiEffrosLifting}; fix one such lift.
Fix a countable
rational $*$-subalgebra dense in $S$, together with countable
dense local tests in $S(V)$ for a compact-containment basis and
all finite unions of its members.  Include rational linear
combinations, products, adjoints, and positive cuts needed for
the algebraic and ideal tests.  Fix $\eta_j\downarrow0$.  At stage $j$
rescale the first $j$ tests and their relevant products into the unit
ball.  Apply \Cref{cc:finite-correction} with a tolerance small enough
that, after rescaling back, its norm and multiplicativity errors on
these tests are at most $\eta_j$.  Denote the resulting map by $L_j$
and its finite repeat size by $N_j$.  In particular,
$\|L_j(a)\|\leq\|a\|+\eta_j$ for every current input $a$.

Write $h^{\mathrm{pad},j}=h\oplus0^{\oplus N_j}$ in the current
row, and subtract it from $L_j$.  The remaining terms
are its compact correction and the $N_j$ small copies of
$\theta$.  They belong to the appropriate matrix ideal $J(V)$
on every tested input in $S(V)$.  On each fixed input $a$, these row
differences also satisfy
\[
 \|L_j(a)-h^{\mathrm{pad},j}(a)\|
 \leq2\|a\|+\eta_j
\]
for all sufficiently large $j$.  Thus they meet the boundedness
hypothesis for prescribed errors in \Cref{seq:diagonal}; no uniform
bound on the global norms of the maps $L_j$ is needed.
For each fixed positive cut
tolerance their row estimates therefore have the form
\begin{equation}\label{cc:row-smallness}
 d_\tau\bigl((|z_{j,n}|-\eps)_+\bigr)
 \le r_{j,n}\,d_\tau(c_j),\qquad
 r_{j,n}\longrightarrow0\quad(n\to\infty),
\end{equation}
where $c_j$ is compactly Cuntz supported in $Q(V)$ and the
bound is uniform in $\tau$.

For each scheduled pair $V\Subset U$, first fix a cut-positive
$b_{V,U}\in Q(U)$ whose ideal contains $Q(V)$.  Compact support
gives $[c_j]\le M_j[b_{V,U}]$ for a finite integer $M_j$.
At stage $j$, wait in its coordinate row until
$M_jr_{j,n}\le 1/j$ for all the finitely many current tests and
cut tolerances.  Wait also until their coordinate norm,
involution and product estimates agree with the rescaled finite-test
bounds, with an additional error tending to zero.  The integers $M_j$
and $N_j$ are chosen before this waiting time; neither needs a bound independent of $j$.
This is the diagonal procedure of \Cref{seq:diagonal}.

Choose coordinate representatives linearly on each finite
rational span, preserving the representatives already needed
for its displayed algebraic tests.  The preceding eventual
estimates allow successive rows to be joined.  Matrix folding
in \Cref{st:folding} is applied only after these choices.  Its
corner identification is adjusted to be within $1/j$ of the
identity on the first $j$ constants of a fixed dense subset of
$Q$.  The resulting corner map on the sequence algebra fixes
$Q$ pointwise.  Apply the corner embeddings to the fixed coordinate
lift of $h$ to define
$h^{\mathrm{pad}}$.

The joined corrected representatives define a contractive,
linear, $*$-preserving, multiplicative map on the dense
rational $*$-subalgebra: each asserted identity and norm bound
is one of the eventually imposed tests.  It extends uniquely
to a homomorphism $\lambda:S\to D$.  The joined estimates
\eqref{cc:row-smallness}, now against the fixed
$b_{V,U}$, put its difference from $h^{\mathrm{pad}}$ in
$J(U)$ on the tests from $S(V)$.  Both resulting maps are
continuous.  Since the ideals $S(V)$ with $V\Subset U$ span
densely in $S(U)$ and $J(U)$ is closed, this proves
\eqref{cc:corrected-difference} for every open $U$.

For completeness, if $\iota_n:Q\to Q$ are the coordinate
corner embeddings, they induce a homomorphism $\iota:D\to D$
with $\iota|_Q=\id_Q$ and $h^{\mathrm{pad}}=\iota h$.
The pullback of a regularized trace along $\iota$ is the
regularized trace of the coordinate weights
$\tau_n\iota_n$, with the same restriction $\sigma$ to $Q$.
Thus every prescribed formula for $h$ of this kind is
preserved.  Applying $\iota$ to an ideal comparison involving
constants also preserves that comparison and those constants.
\end{proof}

\subsection{Recovery of the full trace formula}

We record the trace determination used at the last step.  A
positive element compactly Cuntz supported in the densely
finite ideal of a lower-semicontinuous trace has finite rank.
Indeed that ideal is generated by finite-trace positives;
their positive cuts have finite rank.  Compact containment
dominates the given element by finitely many such cuts.
Conversely, if an ideal is generated by finite-rank positives,
the trace is densely finite on it, since their hereditary
algebras consist of finite-trace elements.

\begin{lemma}[Determination on a finite ideal]\label{cc:finite-trace-determination}
Let $V\subset X$ be open.  Suppose two lower-semicontinuous
traces on $S(V)$ agree on all positive even integer powers of the rich
elementary positive test supply of \Cref{tc:rough-model}, and
these values are finite.  If the cuts of that supply generate
$S(V)$, the two traces are densely finite and equal on $S(V)$.
\end{lemma}

\begin{proof}
Write $\omega_1,\omega_2$ for the traces, and $\omega$ for
either one when making an estimate.
Finite second moments give finite ranks on all positive cuts
of a test $s$: for $\eps>0$,
\[
 d_\omega((s-\eps)_+)\le\eps^{-2}\omega(s^2).
\]
The preceding observation proves dense finiteness for both
traces.  Fix $s$ and a continuous function $g$ on its squared
spectrum which vanishes near zero.  Approximate $g(t)/t$
uniformly by polynomials.  Equality of the moments
$\omega(s^{2k})$ therefore gives equality on $g(s^2)$:
the error is bounded by the uniform polynomial error times
$\omega(s^2)$.  Increasing such functions to spectral support
indicators gives equality of the ranks of every positive cut
of $s^2$.

The test supply contains cuts from dense sets of the two
factors, so these squares approximate every positive
elementary tensor in $S(V)$.  The comparison
\[
 \|a-b\|<\delta
 \quad\Longrightarrow\quad
 (a-t-\delta)_+\precsim(b-t)_+
 \qquad(t\ge0)
\]
transfers equality of all cut ranks to each such elementary
limit, by allowing the cut slack to decrease to zero.  Spectral
integration now gives equality of trace values on all positive
elementary tensors.

To pass from elementary tensors to the whole algebra, let $L$ be a
relatively compact open subset of $(0,1]$ and let
$H_c=\Her(c)\subset P(V)$, where $c$ is compactly Cuntz supported
in $P(V)$.  Choose $f_0\in C_c((0,1])_+$ equal to one on
$\overline L$, and write $c\precsim(b-\delta)_+$ with
$b\in P(V)_+$ and $\delta>0$.  Then
\[
 f_0\otimes c\precsim(f_1\otimes b-\delta/2)_+
\]
for a compactly supported $f_1$ equal to one on $\supp(f_0)$.
The larger tensor belongs to the common finite ideal $S(V)$ already
established above.  Thus $f_0\otimes c$ has finite rank for both
traces.  The support of $C_0(L)\otimes H_c$ is below its support,
so both traces are bounded on this localized algebra, with norms
bounded by that rank.  A positive continuous $P$-valued
function there is uniformly approximated by finite positive
elementary sums, using a scalar partition of unity.
Additivity and boundedness yield equality on this localized
algebra.  Finally take increasing scalar cutoffs and an
increasing cut-supported approximate unit of $P(V)$.
Traciality writes the corresponding compressions as
increasing positive approximations under the square root of
the element being tested.  Lower semicontinuity extends the
equality to every positive element of $S(V)$.
\end{proof}

\begin{theorem}[Existence of cone models]\label{cc:cone-existence}
Let $m$ be a finite full-support Radon measure on $(0,1]$.
There is a homomorphism $p:S\to D$ such that, for every
regularized limit trace $\rho$ with $\sigma=\rho|_Q$,
\begin{equation}\label{cc:full-cone-formula}
 \rho\circ p=m\otimes F(\sigma).
\end{equation}
This includes weights whose densely finite ideal is proper
or zero.
\end{theorem}

\begin{proof}
We first preserve the lower ideal signal through correction.  For a fixed
weight we will then identify its trace on the prescribed finite ideal,
and use that signal to rule out any larger finite ideal.

\smallskip\noindent\emph{Preserve the lower signal.} 
Take $h$ from \Cref{tc:rough-model} and correct it by
\Cref{cc:correction}.  For each controlled positive test
$s=f\otimes a$, put
$\delta(s)=\lambda(s)-h^{\mathrm{pad}}(s)$.
It belongs to $J(U_s)$.  Apply the cone version of
\Cref{sm:stabilizers} once more to obtain a locally small
homomorphism $\theta_1:S\to D$ with
\begin{equation}\label{cc:second-stabilizer}
 \delta(s)\in\Ideal_D(\theta_1(s))
\end{equation}
for every such test.  The countable coefficient form permits
all the required cuts and localized families simultaneously.
Set $p=\lambda\oplus\theta_1$ and fold the two corners,
again fixing the constants.  In the following comparisons
we retain matrix notation for this folding.

The ideal of $p(s)$ contains both $\lambda(s)$ and
$\theta_1(s)$, hence contains $h^{\mathrm{pad}}(s)$ by
\eqref{cc:second-stabilizer}.  The lower signal of
\Cref{tc:rough-model} and the preservation of constants give
\begin{equation}\label{cc:retained-signal}
 Q(U_a)\subset\Ideal_D(p(f\otimes a))
 \qquad(f\ne0)
\end{equation}
for each controlled test.  Also
\begin{equation}\label{cc:total-smallness}
 (p-h^{\mathrm{pad}})(S(U))\subset J(U),
\end{equation}
where zero padding of $h^{\mathrm{pad}}$ is understood.

\smallskip\noindent\emph{Identify the trace on its prescribed finite ideal.} 
Fix a regularized trace $\rho$ and write
$Q(V)=\Fin(\sigma)$.  By \Cref{st:ideals},
$\Fin(F\sigma)=P(V)$.  Put
\[
 \omega=\rho\circ p,\qquad \nu=m\otimes F\sigma.
\]
The trace $\rho$ vanishes on $J(V)$ by
\Cref{seq:smallideal-properties}.  It therefore factors
through the quotient by this zero ideal.  For each controlled
$s=f\otimes a$ localized with its controller in $V$,
\Cref{tc:rough-model} and \eqref{cc:total-smallness} give
\begin{equation}\label{cc:square-agreement}
 \omega(s^2)
 =\rho\bigl(h^{\mathrm{pad}}(s)^2\bigr)
 =\int f^2\,dm\, F\sigma(a^2)
 =\nu(s^2)<\infty.
\end{equation}
Apply the same formula to the positive powers in the supply.
The cuts of these localized elementary tests generate
$S(V)$, including when $V$ is obtained as an increasing union
of the basis opens.  Thus \Cref{cc:finite-trace-determination}
proves
\begin{equation}\label{cc:finite-region}
 S(V)\subset\Fin(\omega),\qquad
 \omega|_{S(V)}=\nu|_{S(V)}.
\end{equation}

\smallskip\noindent\emph{Exclude an additional finite region.} 
We show that the first inclusion is equality.  Suppose the
open set corresponding to $\Fin(\omega)$ contains a
rectangle $L\times U$ not contained in $(0,1]\times V$.
Use a scalar test with compact support in $L$ and a cut
positive $a$ from the basis supply in $P(U)$.  The tests may
be chosen with slightly enlarged supports still in this
rectangle.  Their tensor products are then compactly Cuntz
supported in $\Fin(\omega)$, and consequently have finite
$\omega$-rank.  For any such nonzero scalar test, write
$s=f\otimes a$.  If $c\in Q(U_a)_+$ and $\eps>0$,
\eqref{cc:retained-signal} gives, for suitable finite $M$
and $\eta>0$,
\begin{equation}\label{cc:constant-finiteness}
 [(c-\eps)_+]\le M[(p(s)-\eta)_+]
 \quad\hbox{in }\Cu(D).
\end{equation}
It follows that
\[
 d_\sigma((c-\eps)_+)
 =d_\rho((c-\eps)_+)
 \le M d_\omega(s)<\infty.
\]
Hence $\sigma$ is densely finite on $Q(U_a)$, so
$U_a\subset V$.  The supports $U_a$ of these cut tests
exhaust $U$.  This implies $U\subset V$, contradicting the
choice of the rectangle.  Therefore
\begin{equation}\label{cc:exact-finite-ideal}
 \Fin(\omega)=S(V)=\Fin(\nu).
\end{equation}
Both traces take value infinity on positive elements outside
this ideal.  Together with \eqref{cc:finite-region}, this
proves \eqref{cc:full-cone-formula}.  The argument also
applies when $V$ is empty: there is no finite region to
identify, and the retained signal excludes every nonzero
additional finite ideal.
\end{proof}
